\section{Detection Engineering Fundamentals}

Detection engineering focuses on identifying the behaviors, patterns, and signals that indicate malicious activity within an environment. While threat hunting is exploratory and hypothesis-driven, detection engineering is structured and repeatable. The goal is to translate knowledge of adversary behavior into reliable, high-fidelity detections that operate continuously in a SIEM such as Splunk or Microsoft Sentinel. This section introduces the foundational concepts and provides practical SPL and KQL examples that reflect real-world detection workflows.

A good detection begins with understanding the behavior you want to identify. Instead of relying on static indicators such as IP addresses or file hashes, modern detections focus on techniques—how an adversary executes code, moves laterally, steals credentials, or evades defenses. These behaviors tend to be more stable over time and more resilient to simple changes in tooling. For example, while a malware family may change its command-and-control infrastructure frequently, its use of encoded PowerShell commands or suspicious parent-child process chains often remains consistent.

One of the most common detection patterns involves monitoring for suspicious command-line activity. PowerShell is frequently abused by adversaries because it is built into Windows and provides extensive system access. The following example detects encoded PowerShell commands, which are often used to hide malicious intent:

\textbf{SPL Example}
\begin{lstlisting}
index=sysmon EventCode=1 Image="*powershell.exe"
| search CommandLine="*EncodedCommand*"
| table _time, Computer, User, ParentImage, Image, CommandLine
\end{lstlisting}

\textbf{KQL Equivalent}
\begin{lstlisting}
Sysmon
| where EventID == 1 and Image endswith "powershell.exe"
| where CommandLine contains "EncodedCommand"
| project TimeGenerated, Computer, User, ParentImage, Image, CommandLine
\end{lstlisting}

Another foundational detection pattern focuses on parent-child process anomalies. Legitimate processes tend to follow predictable execution chains, while malicious activity often introduces unusual or rare combinations. The following example identifies uncommon parent-child pairs:

\textbf{SPL Example}
\begin{lstlisting}
index=sysmon EventCode=1
| stats count by ParentImage, Image
| where count < 3
\end{lstlisting}

\textbf{KQL Equivalent}
\begin{lstlisting}
Sysmon
| where EventID == 1
| summarize Count=count() by ParentImage, Image
| where Count < 3
\end{lstlisting}

Authentication-based detections are equally important. Excessive failed logons, unusual logon types, or logons from unexpected hosts can indicate credential misuse or lateral movement. The following example identifies accounts with repeated failed logons:

\textbf{SPL Example}
\begin{lstlisting}
index=wineventlog EventCode=4625
| stats count by AccountName, WorkstationName
| where count > 10
\end{lstlisting}

\textbf{KQL Equivalent}
\begin{lstlisting}
SecurityEvent
| where EventID == 4625
| summarize Count=count() by Account, Workstation
| where Count > 10
\end{lstlisting}

Network-based detections help identify command-and-control activity, data exfiltration, or lateral movement. Rare outbound connections are often a strong signal of malicious behavior:

\textbf{SPL Example}
\begin{lstlisting}
index=sysmon EventCode=3
| stats count by DestinationIp, DestinationPort
| where count < 5
\end{lstlisting}

\textbf{KQL Equivalent}
\begin{lstlisting}
Sysmon
| where EventID == 3
| summarize Count=count() by DestinationIp, DestinationPort
| where Count < 5
\end{lstlisting}

Effective detection engineering also requires tuning and validation. A detection should be tested against real telemetry to ensure it produces meaningful results without overwhelming analysts with noise. This often involves adding filters, whitelisting known-good activity, or enriching events with additional context. Over time, detections should be reviewed and updated as adversary behavior evolves or as new telemetry sources become available.

The examples in this section represent the core building blocks of detection engineering: identifying suspicious behavior, correlating related events, and transforming threat intelligence into actionable, automated logic. These fundamentals will be expanded in later sections into full detection patterns, dashboards, and playbooks that reflect real-world adversary techniques.
